\section{Detailed Methods and Protocol History}
\label{app:detailed-methods}

We use \emph{educational personality} operationally for recurring default
behavior and conditional response patterns of a deployed model in educational
interaction. Let $Y_{mqsar}^{(k)}$ denote whether model configuration $m$ expresses
behavior $k$ on source problem $q$, student state $s$, instruction arm $a$, and
independent request $r$. A default profile summarizes behavior under the
canonical tutor role over a specified distribution of educational inputs. A
conditional profile additionally represents how that behavior varies with
student state. We test both descriptions through predictions on separate source
problems. This operationalization does not require a fixed cross-context model
ranking, and it does not treat any identifiable linguistic signature as a
validated personality construct.

The unit $m$ is a deployment configuration, including its requested and returned
model names, provider, observation period, and generation settings. The five
prospective configurations request MiniMax-M3, MiniMax-M2.7, GLM-5.3,
DeepSeek-V4-Pro, and Doubao-Seed-2.0-Lite. In the measurement pilot, the request
alias GLM-5.2 returned GLM-5.3 throughout; those responses are identified by the
returned name and are not pooled with historical GLM-5.2 results. Provider model
names do not attest to immutable weights, so version and temporal uncertainty
remain limitations of deployment-level inference.




\subsection{Archive Exploration and Prospective Separation}

The archive audit distinguishes raw records from successful deduplicated
responses, and distinguishes tutors from solvers and evaluators. Repeated
questions and conversations can occur under different benchmark names and
prompt variants. Consequently, benchmark names and local item identifiers are
not sufficient independence units. Related source material must stay together
when estimating transfer. Analyses informed by previous archive results remain
exploratory; starting a new conversational context does not make them
prospective.

Across the 899,025 variant-preserving text records, a current-adapter role review
identifies 68,799 records in nine tutor-response settings; most other records
come from evaluation or test-answer tasks. Tutor-role membership does not itself
establish neutral prompts or equivalent historical requests. Appendix~A reports
the complete role inventory and a hash-verified inspection of the 122 tutor
prediction files, including their missing generation-provenance fields.

The controlled protocol uses 32 training questions in 30 conservative source
families and 32 confirmation questions in 32 separate source families. Each
partition has eight questions in each of mathematics, science, language, and
logic/data reasoning. Two training question pairs share an averaging or unit-
conversion family and remain grouped in internal validation. The four pilot
templates and their numerical substitutions are excluded. Confirmation tests
new sources within these four domains, rather than transfer to entirely new
academic domains or proof that a model has never encountered the underlying
concepts.

The questions, correct references, erroneous student attempts, and correct
unfinished work were authored by the research agent. Two model reviewers
checked 64 question packages before freezing, yielding 128 complete reviews.
All five Boolean content checks were positive in every review. Twenty-five
reviews nevertheless included issue text, mostly describing the deliberately
incorrect attempt or the deliberately unfinished work; these comments and their
agent adjudication are retained. This is content quality control without new
human annotation, not a human-gold validation claim.

\subsection{Crossed Educational Inputs}

Each main question has two student-work states, an erroneous attempt and correct
but unfinished work, crossed with two explicit affect statements, calm and
frustrated. These are four controlled input conditions. The work-state contrast
bundles correctness and progress; it does not isolate a pure progress or latent
student-ability effect. The affect statement is an input manipulation, not an
observation of a real learner's emotion.

All prospective main conditions provide the tutor with the same correct
instructor reference for that question, specifying that the student has not
seen it. This equalizes the availability of the solution but does not fully
control comprehension, attention, or instruction-following ability. Results
under this controlled reference condition are reported separately from the
historical archive.

The training stage uses the canonical tutor role on all 32 questions. The
confirmation stage first repeats this design on its 32 new questions. Within
each confirmation domain, four questions are selected by a fixed hash for four
additional role arms: two neutral paraphrases, an instruction to prioritize
student reasoning without revealing the final answer, and an instruction to
explain the solution and state the answer. Each selected question receives every
arm and student state. The neutral arms change only the role wording; the two
policy arms explicitly change the requested teaching behavior.

Every model--input cell receives two requests at temperature 0.7 with an
8,192-token output budget. The fixed design specifies 1,280 training responses,
1,280 canonical confirmation responses, and 2,560 additional intervention
responses. It also includes 160 secondary responses to paired complete/missing-
information probes on eight confirmation sources, plus 40 cross-stage responses
to four training-source sentinels. Sentinels and prompt variants do not increase
the count of independent confirmation sources. Request manifests, input hashes,
parameters, timestamps, and failure histories are recorded. Requests are
interleaved within shared provider concurrency limits rather than run as
separate model blocks.

\subsection{Observable Events and Measurement Development}

Three primary events are fixed from the measurement-development pilot.
\emph{Answer revelation} records an explicit final resolution, including an
incorrect resolution or a correct statement that information is insufficient.
It therefore measures disclosure, not correctness. \emph{Reasoning elicitation}
records an invitation for the student to perform unanswered reasoning, including
imperatives and indirect suggestions; rhetorical questions immediately answered
by the tutor do not qualify. \emph{Affect acknowledgement} records explicit
recognition of feeling or difficulty, excluding generic politeness and praise.
These are behavior endpoints, not three assumed independent personality factors.

Five secondary events cover worked explanation, learner choice, epistemic
qualification, requests for missing facts, and unsupported ability attribution.
All are retained regardless of their results. The pilot showed why a high
aggregate agreement rate is insufficient: requesting missing facts had only one
majority-coded positive example among 160 responses, while unsupported ability
attribution had weak positive agreement. Neither is promoted to a primary
endpoint based on subsequent confirmation results.

MiniMax-M3, DeepSeek-V4-Pro, and GLM-5.3 each code every visible tutor response.
Coders receive the student conversation, the reference target, and the answer
with original line identifiers. They do not receive the generating model name,
system-role assignment, repeat number, or partition label. Natural output style
can still reveal identity. Positive codes must reference existing nonempty
source lines; negative codes have no references. Valid references establish
traceability, not semantic truth. We use majority labels when all three coders
are structurally valid, and retain individual and leave-one-coder analyses,
including fractional means for two-coder ties.

The initial exact-quotation protocol failed on naturally formatted answers,
which motivated line references during measurement development. Original
failures remain archived. The finalized pilot has 480 valid codes for 160
unchanged tutor responses after two structure/completion failures were retried
under identical inputs. Formal structure failures are likewise recorded and
retried without selecting by event label or coder agreement.
A training coding interruption was recovered using saved valid outputs and
the existing two-pass repair allowance. This yielded 3,839 of 3,840 valid codes.
One GLM coding request repeatedly exhausted its 16,384-token completion budget
without visible output. Before forecast locking or confirmation generation, a
disclosed amendment allowed this single missing code a 32,768-token ceiling
and at most two attempts, keeping its model, input, rubric, and temperature
fixed. It completed on the first attempt using 3,511 completion tokens; this
does not identify the cause of the preceding failures. The original and
expanded-budget groups were validated separately before their disjoint codes
were combined. All 1,280 training answers consequently have three valid codes.
The interrupted run, failed attempts, and amendment receipt are retained; the
single expanded-budget code is an explicit deviation from the original
uniform-budget protocol.

\subsection{Prediction Before Confirmation Generation}

For each behavior, four nested logistic predictors model the probability of its
occurrence: a subject-by-student-state baseline; that baseline plus model
identity; an additional model-by-subject block; and an additional
model-by-student-state block. New feature blocks are orthogonalized against
existing blocks using training data and standardized under source-equal
weights. This prevents duplicate main-effect encodings from changing their
ridge penalty when interactions are introduced. The intercept is unpenalized.
The penalty is selected from $\{0.1,1,10,100\}$ by source-grouped training-only
validation. Constant training outcomes use the same smoothed constant across
all four predictors.

The fitted models, tuning choices, input hashes, and confirmation probabilities
are saved before any confirmation tutor response is generated. The generation
runner checks this prediction lock. No confirmation label, answer length, or
same-item cross-model score enters a primary predictor. A secondary competing
forecast uses only training-source profiles of log word count and marked-line
fraction, recomputing those profiles within training validation. Such style
competition is descriptive: teaching policy may itself influence length and
formatting, so it is not a causal adjustment for style.

For each of the three primary events, the two primary comparisons are the Brier
improvement from a default model profile over the situation baseline and from a
student-state profile over the model-by-subject baseline. We average losses
within source and then equally across the 32 confirmation sources. The six
one-sided paired source-level tests use Holm correction. We report absolute
Brier losses, source-bootstrap intervals, and all comparison signs. These
intervals condition on the fitted training profiles and a source-sampling
approximation; they do not describe uncertainty over all models or all
educational contexts.

\subsection{Repetition and Prompt Plasticity}

Repetition analyses report positive, negative, and total event agreement and
profile correspondence, retaining source clustering. Constant behavior can
produce perfect repeat agreement and is reported as such. On the 16 fully
crossed intervention sources, we estimate each model's event-rate difference
between an alternative role and the canonical role. We report average shifts
alongside model and student-state heterogeneity.

For neutral wording, the fixed equivalence reference is $\pm0.10$ in event
probability. We retain the nominal Bonferroni paired-$t$ intervals across the
30 model--event--neutral-arm comparisons. A nonsignificant difference does not
establish equivalence, and constant source differences receive a degeneracy
flag. Further synthetic diagnostics during training show undercoverage in a
sparse boundary case even with nonzero source variance. Before confirmation
generation, we therefore additionally require a bounded-source interval to
support any population-equivalence statement. Scaling Hoeffding's bound to
source differences in $[-1,1]$ and allocating error across both tails and all
30 comparisons gives half-width
$\sqrt{2\log(2\cdot30/0.05)/n}$ \citep{hoeffding1963}. At $n=16$ this is
approximately $0.941$, too wide to support the chosen tolerance. This
conservative sensitivity still assumes independent source observations and
does not make the authored bank representative. We can describe observed
neutral shifts and nominal intervals, while withholding population equivalence.
These disclosed interpretation additions preserve the frozen numerical output
and the six primary prediction tests. The
explicit ASK and EXPLAIN arms quantify policy displacement rather than better
teaching. Finally, answer revelation and reasoning elicitation are analyzed
jointly as four possible combinations. Their evidence lines do not identify
which action occurred first, and we do not infer action order from them.

\subsection{Sensitivity Analyses and Their Timing}

During training coding, after measurement development but before confirmation
responses exist, we additionally fix eight alternative coding panels: the
three-coder mean, each individual coder, each two-coder mean, and the mean after
excluding the generator's model family. For each panel, forecasts are refitted
using its training labels and locked before confirmation generation. They are
evaluated against the corresponding confirmation labels; two-coder ties remain
0.5. This checks dependence on coding choices but cannot remove biases shared
by all coders or supply an external ground truth.

Five further forecast sets each remove one generator configuration from both
training and the confirmation target, refitting the baselines and tuning within
the retained training data. This checks dependence on a single configuration;
it changes the target model mixture and does not test prediction of an unseen
model. Neither sensitivity family replaces the six primary comparisons.

We also resample the 30 training source families 200 times, preserving the
multiplicity of repeated source draws. We retain the selected regularization
relative to total loss weight, without repeating hyperparameter selection, and
recompute training-style profiles in each resample. The resulting gain quantiles
describe training-source sensitivity conditional on the confirmation material
and selected regularization. They are separate from the primary confirmation-
source intervals. These extensions were specified before confirmation outputs,
not preregistered before the entire research project began.

\subsection{Mathematics and Content-Error Sensitivity}

The protocol also requires a mathematics subset and a check of dependence on
obvious content errors. Its implementation is fixed during training coding,
before confirmation generation. We evaluate the existing forecasts on all eight
mathematics sources, covering 320 canonical responses, without refitting.
Separately, MiniMax-M3 and DeepSeek-V4-Pro review all 1,280 canonical responses
for checkable content errors using the student input, reference, and original
answer lines. They distinguish a clear error, no clear error identified, and
uncertainty. Withholding a solution, asking a question, or correctly quoting
and correcting an error is not itself an error. Positive error judgments
require a category and source-line evidence.

Twelve agent-authored controls accompany the responses, giving 2,584 planned
secondary review requests. Each reviewer must match at least five of six
controls in each of the clear-error and no-clear-error classes to meet the
limited diagnostic rule. If either reviewer fails, screened results remain
measurement-limited diagnostics. These controls do not establish correctness
on natural responses, even when the rule is met.

We retain the full primary results and report three additional screens:
excluding unanimously identified errors; excluding errors identified by either
reviewer; and excluding either errors or uncertainty. A flagged response
removes its entire five-model source--student-state--repeat slot so that model
comparisons remain paired. We report retained coverage and source-equal losses
with descriptive source intervals, without new hypothesis tests. These screens
condition on generated responses and change the target distribution; they are
not causal controls for ability. The content review covers the canonical panel,
not all prompt-intervention responses.
